\section{Site selection for omitted moderators}
\label{sec:omitted-moderator}

The fixed-effects response model is a statement about which site-level
variables modify the response. It can fail because a relevant moderator is
omitted from $C$. Omitted-variable bias is especially useful here because it
separates two quantities that are often conflated: the strength of the omitted
moderator in the treatment effect and the extent to which the moderator is
imbalanced between the source sites and the target.

Let $U\in\R^q$ be a site-level moderator with profile $u_s$ at site $s$.
For this section, suppose the treatment-effect contrast obeys the augmented
response model
\begin{equation}
  \E_\sP[Y(1)-Y(0)\mid X,C,U]
  =g_\Delta(X)+C^\top\theta+U^\top\gamma,
  \qquad
  \theta=\beta_1-\beta_0,
  \label{eq:omitted-response-model}
\end{equation}
where $\gamma\in\R^q$ is the effect of the omitted moderator. The analysis in
the rest of the paper fits the model without $U$. Define
\begin{equation}
  \muU(x)=\E_\sP[U\mid X=x],
  \qquad
  \CUcov=\E_\sP[\cov(C,U\mid X)],
  \qquad
  \ellU=\uT-\E_{X\sim\tP}[\muU(X)].
  \label{eq:omitted-objects}
\end{equation}
The target now has an omitted-moderator profile $\uT$. These quantities need
not be observed; they index the sensitivity analysis.

\begin{proposition}[Omitted-moderator bias]
\label{prop:omitted-bias}
Let $\theta^C$ solve the normal equations from fitting the response model
without $U$, so that
\begin{equation}
  \CovC\theta^C=\dmom{1}-\dmom{0}.
  \label{eq:omitted-normal-equation}
\end{equation}
Under \eqref{eq:omitted-response-model},
\begin{equation}
  \dmom{1}-\dmom{0}=\CovC\theta+\CUcov\gamma.
  \label{eq:omitted-moment-shift}
\end{equation}
If $\ct\in\affhull(\pool)$, then the target value reported by the
misspecified model, $\psi_T^C=\bT+\ellvec^\top\theta^C$, differs from the
true target effect by
\begin{equation}
  \psi_T-\psi_T^C
  =\bigl(\ellU-\ellvec^\top\CovC^\dagger\CUcov\bigr)^\top\gamma.
  \label{eq:omitted-bias}
\end{equation}
\end{proposition}

The two terms in \eqref{eq:omitted-bias} have a direct interpretation. The
loading $\ellU$ is the source--target difference in the omitted moderator,
while $\ellvec^\top\CovC^\dagger\CUcov$ is the component that the fitted
context coefficients absorb because $U$ is associated with $C$ in the source
pool. The formula is invariant to the choice of solution $\theta^C$ because
$\ellvec\in\range(\CovC)$ when the target is in the affine hull.

The same decomposition gives a new role to sites that are interior in the
observed context geometry. Let a candidate satisfy
$\cand{k}=\sum_s w_s c_s$ with $\sum_s w_s=1$, and define its omitted-profile
contrast
\begin{equation}
  \udisc{k}=u_k-\sum_s w_su_s.
  \label{eq:omitted-direction}
\end{equation}
The weights may be negative, as they are for affine, rather than convex,
combinations. If $\udisc{k}\neq0$, the candidate is outside the affine hull
of the existing profiles in the augmented $(C,U)$ space even though it is
inside the affine hull in $C$ alone. It therefore tests a direction that the
maintained model treats as irrelevant.

\begin{theorem}[Refinement under an omitted moderator]
\label{thm:omitted-refinement}
Suppose the augmented response model \eqref{eq:omitted-response-model} holds
for the source pool and candidate $k$. Let $\rho_k$ denote the population
contrast between the candidate's conditional treatment effect and the same
affine combination of source-site conditional treatment effects. Then, for
every $x$,
\begin{equation}
\begin{split}
\rho_k
&=\E[Y(1)-Y(0)\mid X=x,C=\cand{k},U=u_k]\\
&\quad-\sum_s w_s\E[Y(1)-Y(0)\mid X=x,C=c_s,U=u_s]
=\udisc{k}^\top\gamma.
\end{split}
\label{eq:omitted-contrast}
\end{equation}
Consider an omitted-moderator sensitivity set
\begin{equation}
  \ovbset
  =\bigl\{(\theta,\gamma):\CovC\theta+\CUcov\gamma
  =\dmom{1}-\dmom{0},\ (\theta,\gamma)\in\restrset_U\bigr\},
  \label{eq:omitted-feasible-set}
\end{equation}
and the target functional
\begin{equation}
  \psi_T(\theta,\gamma)
  =\bT+\ellvec^\top\theta+\ellU^\top\gamma.
  \label{eq:omitted-target-functional}
\end{equation}
After the candidate experiment is observed, the feasible set is
\begin{equation}
  \ovbset^{(k)}
  =\ovbset\cap\bigl\{(\theta,\gamma):\udisc{k}^\top\gamma=\rho_k\bigr\}.
  \label{eq:omitted-feasible-update}
\end{equation}
If the true $(\theta,\gamma)$ belongs to $\ovbset$, then the corresponding
identified-set width satisfies, for every target,
\begin{equation}
  \ovbwidth(\ct,\uT;\pool\cup\{k\})
  \leq \ovbwidth(\ct,\uT;\pool).
  \label{eq:omitted-monotonicity}
\end{equation}
If $\udisc{k}=0$, then $\rho_k=0$ and the candidate adds no omitted-moderator
restriction, so the two widths are equal. If $\udisc{k}\neq0$, the candidate
can strictly reduce the width whenever the restriction
$\udisc{k}^\top\gamma=\rho_k$ cuts a target-relevant direction in $\ovbset$.
\end{theorem}

The theorem is the omitted-variable analogue of \cref{thm:monotone}. Under the
maintained model, a candidate outside the observed affine hull is valuable
because it identifies a new context direction. Under the augmented model, an
interior candidate is valuable when it is outside the affine hull after the
omitted moderator is included. The design criterion is therefore to seek
sites for which $\udisc{k}$ is large in directions of $\gamma$ that matter for
$\ellU$. If $U$ is not measured, this direction must be supplied by a proxy,
substantive prior, or sensitivity range. Without such information, an
unrestricted omitted moderator remains unconstrained by any site-selection
rule based on $C$ alone.
